\section{Requirement Fulfilment}
\label{sec:fulfilment}

Each of the eight requirement families is reported below with what was built
and what evidence exists that it works. A requirement is recorded as met only
where a campaign artefact demonstrates it, not where the implementing code
exists. One is recorded as built but unexercised, and the missing input is
named.

\begin{table}[H]
  \centering\small
  \setlength{\tabcolsep}{5pt}
  \begin{tabularx}{\textwidth}{@{}>{\hsize=0.42\hsize}Y>{\hsize=0.78\hsize}Y>{\hsize=0.38\hsize}Y>{\hsize=1.42\hsize}Y@{}}
    \toprule
    Requirement & Subject & Status & Evidence \\
    \midrule
    ETA-REQ-301 & Defect history as a knowledge source & Built, unexercised & Defect nodes, the ingestion endpoint, defect-weighted retrieval and a defect-history prompt block are live. No campaign exercised them, because no defect export was available \\
    \addlinespace
    ETA-REQ-302 & Execution-path memory as a navigation tree & Met & Median device steps per test fell from 50 to 9. Navigation memory held 123 nodes and 95 recorded dead ends at the end of the sequence \\
    \addlinespace
    ETA-REQ-303 & Incremental experiential learning & Met & Execution logs, mined error patterns and a strategy memory populated from runs and read by the planner \\
    \addlinespace
    ETA-REQ-304 & Multi-dimensional partitioning by profile, platform and application & Met & Five targets across two platforms held in one graph without cross-contamination \\
    \addlinespace
    ETA-REQ-305 & Self-healing and adaptive execution & Met & Failure classification with per-category retry; the livelock detector and the device reset belong to this family \\
    \addlinespace
    ETA-REQ-306 & Regression risk scoring and risk-weighted prioritisation & Met & Per-area risk scores drive prioritisation and are exposed through the risk API and the dashboard \\
    \addlinespace
    ETA-REQ-307 & Test-case quality metrics and feedback loop & Met & Per-test effectiveness recorded and fed back into generation; duplicate suppression operates against executed-test titles \\
    \addlinespace
    ETA-REQ-308 & Anomaly detection over execution patterns & Met & Failure-rate spikes, duration regressions and new-failure detection implemented and surfaced \\
    \bottomrule
  \end{tabularx}
\end{table}

\subsection{On ETA-REQ-301}

This is the one family that has not been exercised against real data, and the
distinction matters enough to state rather than bury in a table cell.

Everything specified is built. Defect nodes carry severity, area and
root-cause category. The ingestion endpoint accepts a defect export. Defect
density weights the retrieval score, and a defect-history block is injected
into generation with its own character budget. The path has been exercised on
synthetic data and behaves as specified.

What has not happened is a campaign in which real defect history changed which
tests were chosen. That requires a defect export from a real product, which is
the first of the two requests in \Cref{sec:next}.

\subsection{Non-functional requirements}

ETA-NFR-001 requires retrieval endpoints to respond within 500\,ms at the 95th
percentile on a graph of up to 100{,}000 nodes. ETA-NFR-002 requires support for
projects with up to 10{,}000 test cases, 5{,}000 defects and 50{,}000 navigation
nodes without degradation.

Both are met at the scale reached so far, and neither has been tested at the
stated ceiling. The largest graph built during this work is roughly three
orders of magnitude below the ETA-NFR-002 limits, so what can honestly be
reported is that no degradation was observed in the range exercised, not that
the ceiling has been demonstrated. Load testing at the specified scale is
straightforward and has not been done.

\subsection{Implementation phasing}

ETA-TR-2026-001 proposes six phases totalling 17 to 23 weeks. The work was
carried out as ten numbered packages, each with a verification gate, built and
integrated in sequence between July and August 2026. The mapping is
one-to-many rather than one-to-one: the phases describe capability groupings
and the packages describe deliverable units, and several packages cut across
phases because the graph schema had to be settled before the features that
depend on it could be built independently.
